% =========================================================
% Chapter: Magnetic Fields Produced by Electric Currents
% POSN 1 Physics Preparation Notes
% Language: Thai
% Compiler: XeLaTeX
%
% Required environments/packages from the main project:
%   graphicx, siunitx, tcolorbox
%   examplebox, exercisebox
% =========================================================

\chapter{สนามแม่เหล็กจากวงจรไฟฟ้า}

\begin{center}
    {\Large \textbf{Sources of the Magnetic Field}}\\[4pt]
    {\normalsize เอกสารเตรียมสอบคัดเลือก สอวน. ฟิสิกส์ ค่าย 1}\\[6pt]
    {\normalsize จัดทำโดย \textbf{Tames Thanakrit Damduan}}\\[2pt]
    {\footnotesize \textcopyright\ 2026 Tames Thanakrit Damduan. All rights reserved.}\\[8pt]
\end{center}

\section*{เป้าหมายของบทเรียน}

หลังจากเรียนบทนี้ นักเรียนควรทำสิ่งต่อไปนี้ได้

\begin{enumerate}
    \item อธิบายได้ว่าประจุไฟฟ้าที่เคลื่อนที่หรือกระแสไฟฟ้าเป็นแหล่งกำเนิดสนามแม่เหล็ก
    \item ใช้กฎมือขวาหาทิศของสนามแม่เหล็กจากลวดที่มีกระแสไหลได้
    \item ใช้กฎของบิโอต์--ซาวาร์หาสนามจากลวดตรง ส่วนโค้งวงกลม และวงลวดได้
    \item อนุมานสนามที่ศูนย์กลางและบนแกนของวงลวดจากกฎของบิโอต์--ซาวาร์ได้
    \item ใช้หลักการซ้อนทับรวมสนามจากลวดหลายส่วนได้อย่างถูกต้อง
    \item อนุมานแรงต่อหนึ่งหน่วยความยาวระหว่างลวดตัวนำขนานสองเส้นได้
    \item ใช้กฎของแอมแปร์กับระบบที่มีสมมาตรสูงได้
    \item อนุมานสนามภายในและภายนอกลวดทรงกระบอก โซลินอยด์ และโทรอยด์ได้
    \item อธิบายกฎของเกาส์สำหรับสนามแม่เหล็กและความหมายของเส้นสนามปิดได้
    \item แก้ข้อสอบจริง สอวน. ที่เกี่ยวข้องกับสนามแม่เหล็กจากกระแสไฟฟ้าได้
\end{enumerate}

\begin{tcolorbox}[title={แนวคิดสำคัญของบทนี้}]
ในบทวงจรไฟฟ้า เราศึกษาว่าประจุเคลื่อนที่อย่างไรในตัวนำ บทนี้ศึกษาผลที่เกิดขึ้นรอบตัวนำเมื่อประจุเหล่านั้นเคลื่อนที่ กล่าวคือ กระแสไฟฟ้าสร้างสนามแม่เหล็ก

เครื่องมือหลักมีสองชนิด

\[
\text{กฎของบิโอต์--ซาวาร์: ใช้ได้กับรูปร่างของกระแสทั่วไป}
\]

\[
\text{กฎของแอมแปร์: ใช้ง่ายเมื่อระบบมีสมมาตรสูง}
\]
\end{tcolorbox}

% =========================================================
\section{จากกระแสไฟฟ้าสู่สนามแม่เหล็ก}

\subsection{กระแสไฟฟ้าคือประจุที่เคลื่อนที่}

จากบทวงจรไฟฟ้า กระแสไฟฟ้าถูกนิยามว่าเป็นอัตราการไหลของประจุผ่านหน้าตัดของตัวนำ

\begin{tcolorbox}[title={นิยาม: กระแสไฟฟ้า}]
\[
\boxed{I\equiv\frac{dq}{dt}}
\]
\end{tcolorbox}

เมื่อประจุไฟฟ้าเคลื่อนที่ ประจุนั้นสร้างสนามแม่เหล็กในบริเวณรอบตัวมัน ดังนั้นลวดที่มีกระแสไหลจึงเป็นแหล่งกำเนิดสนามแม่เหล็ก

ในบทนี้ คำว่า ``สนามแม่เหล็ก'' หมายถึงเวกเตอร์

\[
\vec{B}
\]

ซึ่งมักเรียกอย่างเต็มว่า \textbf{ความหนาแน่นฟลักซ์แม่เหล็ก} และมีหน่วยเทสลา

\[
\si{T}
\]

\subsection{นิยามเชิงปฏิบัติการของสนามแม่เหล็ก}

เมื่อประจุ $q$ เคลื่อนที่ด้วยความเร็ว $\vec{v}$ ในสนามแม่เหล็ก $\vec{B}$ แรงแม่เหล็กเขียนได้เป็น

\begin{tcolorbox}[title={กฎแรงแม่เหล็กและนิยามเชิงปฏิบัติการของ $\vec{B}$}]
\[
\boxed{\vec{F}_B=q\vec{v}\times\vec{B}}
\]

ขนาดของแรงคือ

\[
\boxed{F_B=|q|vB\sin\theta}
\]

โดย $\theta$ เป็นมุมระหว่าง $\vec{v}$ กับ $\vec{B}$
\end{tcolorbox}

ถ้า $\vec{v}\perp\vec{B}$ จะได้

\[
F_B=|q|vB
\]

จึงสามารถวัดหรือกำหนดขนาดของสนามแม่เหล็กได้จาก

\[
\boxed{B\equiv\frac{F_B}{|q|v}}
\qquad
(\vec{v}\perp\vec{B})
\]

หน่วยของสนามแม่เหล็กจึงเป็น

\[
1\,\si{T}
=
1\,\si{N/(C.m.s^{-1})}
=
1\,\si{N/(A.m)}
\]

ทิศของแรงบนประจุบวกหาได้จากกฎมือขวาของผลคูณเวกเตอร์ ส่วนแรงบนประจุลบมีทิศตรงข้าม

\begin{tcolorbox}[title={ผลสำคัญ}]
เพราะ $\vec{F}_B$ ตั้งฉากกับ $\vec{v}$ เสมอ

\[
\vec{F}_B\cdot\vec{v}=0
\]

สนามแม่เหล็กจึงไม่ทำงานต่อประจุจุดโดยตรง สนามแม่เหล็กเปลี่ยนทิศของความเร็วได้ แต่ไม่เปลี่ยนอัตราเร็วของประจุเพียงลำพัง
\end{tcolorbox}

\subsection{อนุมานแรงแม่เหล็กบนส่วนเล็กของลวด}

พิจารณาลวดพื้นที่หน้าตัด $A$ มีพาหะประจุจำนวน $n$ ตัวต่อหนึ่งหน่วยปริมาตร แต่ละตัวมีประจุ $q$ และมีความเร็วดริฟต์ $\vec{v}_d$

ในส่วนลวดยาว $d\ell$ จำนวนพาหะประจุคือ

\[
dN=nA\,d\ell
\]

แรงแม่เหล็กรวมบนพาหะทั้งหมดในส่วนลวดนี้คือ

\[
d\vec{F}
=
(dN)q\vec{v}_d\times\vec{B}
\]

แทน $dN=nA\,d\ell$

\[
d\vec{F}
=
nA\,d\ell\,q\vec{v}_d\times\vec{B}
\]

จากบทวงจรไฟฟ้า

\[
I=nqAv_d
\]

และให้เวกเตอร์ $d\vec{\ell}$ ชี้ตามทิศกระแสตามแบบ จึงได้

\begin{tcolorbox}[title={แรงแม่เหล็กบนส่วนกระแส ซึ่งอนุมานจากแรงบนประจุ}]
\[
\boxed{d\vec{F}=I\,d\vec{\ell}\times\vec{B}}
\]
\end{tcolorbox}

ถ้าลวดตรงยาว $L$ อยู่ในสนามสม่ำเสมอ

\[
\boxed{\vec{F}=I\vec{L}\times\vec{B}}
\]

และ

\[
\boxed{F=ILB\sin\theta}
\]

สมการนี้จะถูกใช้ภายหลังเพื่ออนุมานแรงระหว่างลวดที่มีกระแสไหล

% =========================================================
\section{กฎของบิโอต์--ซาวาร์}

\subsection{กฎเชิงทดลองสำหรับสนามจากส่วนกระแส}

ให้ส่วนเล็กของลวดมีเวกเตอร์ความยาว $d\vec{\ell}$ ชี้ตามทิศกระแส $I$ และให้จุดสังเกตอยู่ห่างจากส่วนกระแสเป็นระยะ $r$ โดย $\hat{r}$ ชี้จากส่วนกระแสไปยังจุดสังเกต

\begin{tcolorbox}[title={กฎของบิโอต์--ซาวาร์}]
\[
\boxed{
 d\vec{B}
 =
 \frac{\mu_0}{4\pi}
 \frac{I\,d\vec{\ell}\times\hat{r}}{r^2}}
\]

ในรูปขนาด

\[
\boxed{
 dB
 =
 \frac{\mu_0}{4\pi}
 \frac{I\,d\ell\sin\theta}{r^2}}
\]
\end{tcolorbox}

สมการนี้เป็นกฎที่สรุปจากการทดลอง ไม่ใช่นิยามของสนามแม่เหล็ก

ค่าคงที่ $\mu_0$ เรียกว่า \textbf{สภาพซึมผ่านได้ทางแม่เหล็กของสุญญากาศ} ในโจทย์ สอวน. ใช้

\[
\boxed{
\mu_0=4\pi\times10^{-7}\,\si{T.m/A}}
\]

จากสมการเห็นว่า

\begin{enumerate}
    \item $dB\propto I$
    \item $dB\propto d\ell$
    \item $dB\propto \sin\theta$
    \item $dB\propto 1/r^2$ สำหรับส่วนกระแสเล็กหนึ่งส่วน
\end{enumerate}

\subsection{ทิศของสนามจากกฎมือขวา}

ทิศของ $d\vec{B}$ คือทิศของ

\[
d\vec{\ell}\times\hat{r}
\]

ใช้นิ้วชี้มือขวาชี้ตาม $d\vec{\ell}$ งอนิ้วไปทาง $\hat{r}$ นิ้วหัวแม่มือชี้ทิศของ $d\vec{B}$

สำหรับลวดตรงยาว เส้นสนามเป็นวงกลมร่วมศูนย์กลางรอบลวด ใช้กฎกำมือขวาได้ดังนี้

\begin{enumerate}
    \item นิ้วหัวแม่มือชี้ตามทิศกระแส
    \item นิ้วที่โค้งรอบลวดชี้ตามทิศของสนามแม่เหล็ก
\end{enumerate}

\subsection{หลักการซ้อนทับ}

สนามจากลวดทั้งเส้นต้องรวมสนามจากส่วนกระแสทุกส่วนแบบเวกเตอร์

\[
\vec{B}=\int d\vec{B}
\]

แทนกฎของบิโอต์--ซาวาร์

\begin{tcolorbox}[title={สนามจากการกระจายกระแสต่อเนื่อง}]
\[
\boxed{
\vec{B}
=
\frac{\mu_0 I}{4\pi}
\int
\frac{d\vec{\ell}\times\hat{r}}{r^2}}
\]
\end{tcolorbox}

ถ้ามีลวดหลายส่วนหรือหลายวงจร

\[
\boxed{
\vec{B}_{\text{net}}
=
\vec{B}_1+\vec{B}_2+\vec{B}_3+\cdots}
\]

\begin{tcolorbox}[title={ส่วนลวดที่ให้สนามเป็นศูนย์ ณ จุดสังเกต}]
ถ้าแนวของส่วนกระแสชี้ตรงเข้าหาหรือออกจากจุดสังเกต

\[
d\vec{\ell}\parallel\hat{r}
\]

จึงมี

\[
d\vec{\ell}\times\hat{r}=0
\]

ส่วนลวดนั้นไม่สร้างสนามที่จุดดังกล่าว แม้จะมีกระแสไหลอยู่
\end{tcolorbox}

% =========================================================
\section{สนามแม่เหล็กจากลวดรูปทรงสำคัญ}

\subsection{ลวดตรงยาวจำกัด}

วางลวดตามแกน $x$ และให้จุด $P$ อยู่ห่างจากลวดในแนวตั้งฉากเป็นระยะ $a$

เลือกส่วนกระแสยาว $dx$ ที่ตำแหน่ง $x$ ระยะจากส่วนกระแสถึง $P$ คือ

\[
r=\sqrt{x^2+a^2}
\]

จากเรขาคณิต

\[
\sin\theta=\frac{a}{r}
\]

ใช้กฎของบิโอต์--ซาวาร์

\[
dB
=
\frac{\mu_0 I}{4\pi}
\frac{dx\sin\theta}{r^2}
\]

แทน $\sin\theta=a/r$

\[
dB
=
\frac{\mu_0 I}{4\pi}
\frac{a\,dx}{(x^2+a^2)^{3/2}}
\]

ดังนั้น

\[
B
=
\frac{\mu_0 I}{4\pi}
\int_{x_1}^{x_2}
\frac{a\,dx}{(x^2+a^2)^{3/2}}
\]

สังเกตว่า

\[
\frac{d}{dx}
\left(
\frac{x}{a\sqrt{x^2+a^2}}
\right)
=
\frac{a}{(x^2+a^2)^{3/2}}
\]

จึงได้

\[
B
=
\frac{\mu_0 I}{4\pi a}
\left[
\frac{x}{\sqrt{x^2+a^2}}
\right]_{x_1}^{x_2}
\]

ถ้าให้ $\alpha$ และ $\beta$ เป็นมุมจากเส้นตั้งฉากไปยังปลายลวดด้านซ้ายและด้านขวา จะมี

\[
\sin\alpha
=
\frac{|x_1|}{\sqrt{x_1^2+a^2}}
\]

\[
\sin\beta
=
\frac{x_2}{\sqrt{x_2^2+a^2}}
\]

จึงได้

\begin{tcolorbox}[title={สนามจากลวดตรงยาวจำกัด}]
\[
\boxed{
B
=
\frac{\mu_0 I}{4\pi a}
(\sin\alpha+\sin\beta)}
\]
\end{tcolorbox}

ทิศตั้งฉากกับระนาบที่ประกอบด้วยลวดและจุด $P$ และหาได้ด้วยกฎมือขวา

\subsection{ลวดตรงยาวมากและลวดกึ่งอนันต์}

สำหรับลวดตรงยาวมากทั้งสองด้าน

\[
\alpha=\beta=90^\circ
\]

แทนในสมการลวดจำกัด

\[
B
=
\frac{\mu_0 I}{4\pi a}
(1+1)
\]

จึงได้

\begin{tcolorbox}[title={สนามจากลวดตรงยาวมาก}]
\[
\boxed{B=\frac{\mu_0 I}{2\pi a}}
\]
\end{tcolorbox}

สำหรับลวดกึ่งอนันต์ที่ปลายลวดอยู่ตรงจุดเท้าฉาก

\[
\alpha=0,
\qquad
\beta=90^\circ
\]

จึงได้

\begin{tcolorbox}[title={สนามจากลวดตรงกึ่งอนันต์}]
\[
\boxed{B=\frac{\mu_0 I}{4\pi a}}
\]
\end{tcolorbox}

\subsection{สนามที่ศูนย์กลางของส่วนโค้งวงกลม}

ให้ส่วนโค้งวงกลมรัศมี $R$ รองรับมุม $\phi$ เรเดียน และมีกระแส $I$ ไหล

ที่ศูนย์กลาง ทุกส่วนกระแสมี

\[
r=R
\]

และ

\[
d\vec{\ell}\perp\hat{r}
\]

จึงมี $\sin\theta=1$ และ

\[
dB
=
\frac{\mu_0 I}{4\pi}
\frac{d\ell}{R^2}
\]

จากความยาวส่วนโค้ง

\[
d\ell=R\,d\varphi
\]

ดังนั้น

\[
dB
=
\frac{\mu_0 I}{4\pi R}
\,d\varphi
\]

อินทิเกรตตั้งแต่ $0$ ถึง $\phi$

\[
B
=
\frac{\mu_0 I}{4\pi R}
\int_0^\phi d\varphi
\]

\begin{tcolorbox}[title={สนามที่ศูนย์กลางของส่วนโค้งวงกลม}]
\[
\boxed{
B
=
\frac{\mu_0 I\phi}{4\pi R}}
\qquad
(\phi\text{ ต้องอยู่ในหน่วยเรเดียน})
\]
\end{tcolorbox}

กรณีครึ่งวงกลม $\phi=\pi$

\[
\boxed{B_{\text{semi}}=\frac{\mu_0 I}{4R}}
\]

กรณีหนึ่งในสี่วงกลม $\phi=\pi/2$

\[
\boxed{B_{\text{quarter}}=\frac{\mu_0 I}{8R}}
\]

กรณีเต็มวง $\phi=2\pi$

\[
\boxed{B_{\text{loop}}=\frac{\mu_0 I}{2R}}
\]

ถ้ามี $N$ รอบซ้อนกันและแต่ละรอบมีกระแส $I$

\[
\boxed{B_{\text{center}}=\frac{\mu_0 NI}{2R}}
\]

\begin{tcolorbox}[title={ข้อควรระวังเรื่องส่วนลวดแนวรัศมี}]
ถ้าส่วนลวดตรงวางตามแนวรัศมีจากศูนย์กลาง ส่วนลวดนั้นให้สนามเป็นศูนย์ที่ศูนย์กลาง เพราะ

\[
d\vec{\ell}\parallel\hat{r}
\]

ดังนั้นในลวดที่ประกอบด้วยส่วนโค้งและส่วนตรงแนวรัศมี ให้รวมเฉพาะสนามจากส่วนโค้ง
\end{tcolorbox}

\begin{examplebox}{ตัวอย่าง: ส่วนโค้งสองรัศมี}
ลวดประกอบด้วยครึ่งวงกลมรัศมี $R_1$ และ $R_2$ กระแส $I$ ไหลในทิศที่ทำให้สนามจากทั้งสองส่วนชี้ทิศเดียวกัน จงหาสนามที่ศูนย์กลางร่วม

แต่ละครึ่งวงกลมให้สนาม

\[
B_1=\frac{\mu_0 I}{4R_1}
\]

\[
B_2=\frac{\mu_0 I}{4R_2}
\]

เพราะทิศเดียวกันจึงบวกขนาด

\[
\boxed{
B_{\text{net}}
=
\frac{\mu_0 I}{4}
\left(
\frac{1}{R_1}+\frac{1}{R_2}
\right)}
\]

ถ้ากระแสวนตรงข้ามกัน ให้ลบขนาดแทน
\end{examplebox}

\subsection{สนามบนแกนของวงลวดกลม}

ให้วงลวดรัศมี $R$ อยู่ในระนาบหนึ่ง และให้จุด $P$ อยู่บนแกนตั้งฉากผ่านศูนย์กลาง ห่างจากศูนย์กลางเป็นระยะ $x$

ระยะจากทุกส่วนกระแสถึง $P$ เท่ากันคือ

\[
r=\sqrt{R^2+x^2}
\]

เพราะ $d\vec{\ell}$ เป็นแนวสัมผัสวงกลม จึงตั้งฉากกับ $\vec{r}$ และมี

\[
dB
=
\frac{\mu_0 I}{4\pi}
\frac{d\ell}{r^2}
\]

ด้วยสมมาตร องค์ประกอบที่ตั้งฉากกับแกนหักล้างกัน เหลือเฉพาะองค์ประกอบตามแกน

จากเรขาคณิต องค์ประกอบตามแกนของ $d\vec{B}$ มีตัวคูณ

\[
\frac{R}{r}
\]

ดังนั้น

\[
dB_x
=
\frac{\mu_0 I}{4\pi}
\frac{R\,d\ell}{r^3}
\]

รวมรอบวง โดย

\[
\oint d\ell=2\pi R
\]

จะได้

\[
B_x
=
\frac{\mu_0 I}{4\pi}
\frac{R}{r^3}
(2\pi R)
\]

แทน $r^2=R^2+x^2$

\begin{tcolorbox}[title={สนามบนแกนของวงลวดหนึ่งรอบ}]
\[
\boxed{
B(x)
=
\frac{\mu_0 I R^2}
{2(R^2+x^2)^{3/2}}}
\]
\end{tcolorbox}

ถ้ามี $N$ รอบ

\[
\boxed{
B(x)
=
\frac{\mu_0 NI R^2}
{2(R^2+x^2)^{3/2}}}
\]

ที่ศูนย์กลาง $x=0$

\[
B(0)
=
\frac{\mu_0 NI R^2}{2R^3}
=
\frac{\mu_0 NI}{2R}
\]

จึงสอดคล้องกับผลที่ได้จากส่วนโค้งเต็มวง

ถ้า $x\gg R$

\[
(R^2+x^2)^{3/2}\approx x^3
\]

จึงได้พฤติกรรมระยะไกล

\[
B(x)\approx\frac{\mu_0 NI R^2}{2x^3}
\]

\subsection{วงแหวนประจุที่หมุน}

ให้ประจุรวม $Q$ กระจายสม่ำเสมอบนวงแหวนรัศมี $R$ และหมุนด้วยอัตราเร็วเชิงมุม $\omega$

จากนิยามอัตราเร็วเชิงมุม

\[
\omega=\frac{2\pi}{T}
\]

ดังนั้นคาบคือ

\[
T=\frac{2\pi}{\omega}
\]

ในหนึ่งคาบ ประจุ $Q$ ผ่านจุดใดจุดหนึ่งครบหนึ่งครั้ง จากนิยามกระแส

\[
I=\frac{Q}{T}
\]

แทน $T=2\pi/\omega$

\begin{tcolorbox}[title={กระแสสมมูลของวงแหวนประจุที่หมุน}]
\[
\boxed{I=\frac{Q\omega}{2\pi}}
\]
\end{tcolorbox}

นำไปแทนในสนามที่ศูนย์กลางของวงลวด

\[
B_O=\frac{\mu_0 I}{2R}
\]

จะได้

\begin{tcolorbox}[title={สนามที่ศูนย์กลางของวงแหวนประจุที่หมุน}]
\[
\boxed{
B_O
=
\frac{\mu_0 Q\omega}{4\pi R}}
\]
\end{tcolorbox}

ทิศของสนามหาโดยพิจารณาทิศกระแสตามแบบ ถ้าประจุเป็นบวก ทิศกระแสเหมือนทิศการหมุนของประจุ ถ้าประจุเป็นลบ ทิศกระแสตรงข้ามกับทิศการหมุน

% =========================================================
\section{แรงระหว่างตัวนำที่มีกระแสไหล}

\subsection{แรงระหว่างลวดตรงขนานสองเส้น}

ให้ลวดเส้นที่ 1 และ 2 ยาวมาก วางขนานกัน ห่างกัน $d$ และมีกระแส $I_1$ กับ $I_2$

สนามจากลวดเส้นที่ 1 ณ ตำแหน่งของลวดเส้นที่ 2 คือ

\[
B_1=\frac{\mu_0 I_1}{2\pi d}
\]

แรงบนลวดเส้นที่ 2 ยาว $L$ คือ

\[
F_2=I_2LB_1
\]

เพราะลวดตั้งฉากกับสนามที่ตำแหน่งนั้น

แทน $B_1$

\[
F_2
=
I_2L
\left(
\frac{\mu_0 I_1}{2\pi d}
\right)
\]

หารด้วย $L$

\begin{tcolorbox}[title={แรงต่อหนึ่งหน่วยความยาวระหว่างลวดขนาน}]
\[
\boxed{
\frac{F}{L}
=
\frac{\mu_0 I_1I_2}{2\pi d}}
\]
\end{tcolorbox}

ใช้กฎมือขวาหาทิศได้ว่า

\begin{itemize}
    \item กระแสทิศเดียวกัน: ลวดดึงดูดกัน
    \item กระแสสวนทางกัน: ลวดผลักกัน
\end{itemize}

\subsection{แรงโดยประมาณระหว่างวงลวดสองวงที่อยู่ใกล้กัน}

ให้วงลวดร่วมแกนสองวงมีรัศมี $R$ ห่างกัน $d$ โดย

\[
d\ll R
\]

และมีกระแส $I$ วนทิศเดียวกัน

ในบริเวณเล็ก ๆ ของวงลวด ความโค้งน้อยมากเมื่อเทียบกับระยะ $d$ จึงประมาณส่วนลวดตรงกันข้ามเป็นลวดขนานได้

แรงต่อหนึ่งหน่วยความยาวคือ

\[
\frac{F}{L}
\approx
\frac{\mu_0 I^2}{2\pi d}
\]

ความยาวรอบวงคือ

\[
L=2\pi R
\]

ดังนั้น

\[
F
\approx
\frac{\mu_0 I^2}{2\pi d}
(2\pi R)
\]

\begin{tcolorbox}[title={แรงโดยประมาณระหว่างวงลวดใกล้กัน}]
\[
\boxed{
F\approx\frac{\mu_0 I^2R}{d}}
\qquad
(d\ll R)
\]

ถ้ากระแสวนทิศเดียวกัน แรงเป็นแรงดึงดูด
\end{tcolorbox}

สูตรนี้เป็นการประมาณแบบเฉพาะที่ ไม่ใช่สูตรทั่วไปของวงลวดสองวงเมื่อ $d$ มีขนาดเทียบได้กับ $R$

% =========================================================
\section{กฎของแอมแปร์}

\subsection{การอินทิเกรตตามเส้นปิด}

พิจารณาเส้นทางปิดสมมติ $C$ และส่วนเล็กของเส้นทาง $d\vec{\ell}$

ปริมาณ

\[
\oint_C\vec{B}\cdot d\vec{\ell}
\]

เรียกว่าอินทิกรัลเชิงเส้นของสนามแม่เหล็กรอบเส้นปิด

\begin{tcolorbox}[title={กฎของแอมแปร์สำหรับกระแสคงตัว}]
\[
\boxed{
\oint_C\vec{B}\cdot d\vec{\ell}
=
\mu_0I_{\text{enc}}}
\]
\end{tcolorbox}

$I_{\text{enc}}$ คือผลรวมเชิงเครื่องหมายของกระแสที่ทะลุพื้นผิวซึ่งมีขอบเป็นเส้น $C$

ทิศบวกของกระแสสัมพันธ์กับทิศเดินรอบเส้น $C$ ด้วยกฎมือขวา

\begin{tcolorbox}[title={กฎของแอมแปร์เหมาะกับกรณีใด}]
กฎของแอมแปร์เป็นจริงทั่วไปสำหรับกระแสคงตัว แต่จะใช้หาค่า $B$ ได้ง่ายเมื่อสมมาตรทำให้

\begin{enumerate}
    \item $B$ คงที่บนบางส่วนของเส้นแอมแปร์
    \item $\vec{B}\parallel d\vec{\ell}$ หรือ $\vec{B}\perp d\vec{\ell}$
    \item สามารถเขียนอินทิกรัลเป็น $B$ คูณความยาวได้
\end{enumerate}
\end{tcolorbox}

\subsection{อนุมานสนามของลวดตรงยาวมากด้วยกฎของแอมแปร์}

เลือกเส้นแอมแปร์เป็นวงกลมรัศมี $r$ รอบลวด

จากสมมาตร สนามมีขนาดคงที่บนวงกลมและสัมผัสวงกลมทุกจุด

\[
\vec{B}\cdot d\vec{\ell}=B\,d\ell
\]

ดังนั้น

\[
\oint\vec{B}\cdot d\vec{\ell}
=
B\oint d\ell
\]

\[
\oint d\ell=2\pi r
\]

ใช้กฎของแอมแปร์

\[
B(2\pi r)=\mu_0I
\]

จึงได้

\[
\boxed{B=\frac{\mu_0I}{2\pi r}}
\]

ซึ่งตรงกับผลจากกฎของบิโอต์--ซาวาร์

\subsection{สนามภายในลวดทรงกระบอกที่กระแสกระจายสม่ำเสมอ}

ให้ลวดทรงกระบอกรัศมี $R$ มีกระแสรวม $I$ กระจายสม่ำเสมอทั่วพื้นที่หน้าตัด

\begin{tcolorbox}[title={นิยาม: ความหนาแน่นกระแสสม่ำเสมอ}]
\[
\boxed{J\equiv\frac{I}{\pi R^2}}
\]
\end{tcolorbox}

สำหรับจุดภายในลวดที่ห่างแกน $r<R$ กระแสที่ถูกล้อมด้วยวงแอมแปร์รัศมี $r$ คือ

\[
I_{\text{enc}}
=J\pi r^2
\]

แทน $J=I/(\pi R^2)$

\[
I_{\text{enc}}
=
I\frac{r^2}{R^2}
\]

ใช้กฎของแอมแปร์

\[
B(2\pi r)
=
\mu_0I\frac{r^2}{R^2}
\]

ดังนั้น

\begin{tcolorbox}[title={สนามภายในและภายนอกลวดทรงกระบอก}]
สำหรับ $0\le r<R$

\[
\boxed{
B(r)=\frac{\mu_0Ir}{2\pi R^2}}
\]

สำหรับ $r\ge R$

\[
\boxed{
B(r)=\frac{\mu_0I}{2\pi r}}
\]
\end{tcolorbox}

ภายในลวด

\[
B\propto r
\]

ภายนอกลวด

\[
B\propto\frac{1}{r}
\]

และทั้งสองสมการให้ค่าเดียวกันที่ผิว $r=R$

\subsection{โทรอยด์}

โทรอยด์คือขดลวด $N$ รอบพันรอบแกนรูปวงแหวน

เลือกเส้นแอมแปร์เป็นวงกลมรัศมี $r$ ภายในบริเวณขดลวด สนามสัมผัสวงกลมและมีขนาดคงที่โดยสมมาตร

\[
\oint\vec{B}\cdot d\vec{\ell}
=B(2\pi r)
\]

เส้นแอมแปร์ล้อมกระแสจากขดลวด $N$ รอบ

\[
I_{\text{enc}}=NI
\]

ใช้กฎของแอมแปร์

\[
B(2\pi r)=\mu_0NI
\]

\begin{tcolorbox}[title={สนามภายในโทรอยด์อุดมคติ}]
\[
\boxed{
B(r)=\frac{\mu_0NI}{2\pi r}}
\]
\end{tcolorbox}

สนามภายในโทรอยด์ไม่คงที่อย่างสมบูรณ์ เพราะแปรผกผันกับ $r$ ถ้ารัศมีหลักของโทรอยด์มีขนาดใหญ่มากเมื่อเทียบกับความหนาของหน้าตัด สนามภายในจึงเกือบสม่ำเสมอ

สำหรับโทรอยด์อุดมคติที่พันแน่น สนามในช่องกลางและภายนอกมีค่าน้อยมาก อย่างไรก็ตาม $\oint\vec{B}\cdot d\vec{\ell}=0$ เพียงอย่างเดียวไม่ได้รับประกันว่า $B=0$ ถ้าไม่มีสมมาตรเพียงพอ

\subsection{โซลินอยด์ยาว}

โซลินอยด์คือเส้นลวดยาวพันเป็นเกลียวจำนวนมากรอบแกน

\begin{tcolorbox}[title={นิยาม: จำนวนรอบต่อหนึ่งหน่วยความยาว}]
ถ้าโซลินอยด์ยาว $L$ และมี $N$ รอบ

\[
\boxed{n\equiv\frac{N}{L}}
\]
\end{tcolorbox}

สำหรับโซลินอยด์อุดมคติที่ยาวมาก สนามภายในเกือบสม่ำเสมอและขนานแกน ส่วนสนามภายนอกมีค่าน้อยมาก

เลือกเส้นแอมแปร์รูปสี่เหลี่ยม โดยด้านหนึ่งยาว $\ell$ อยู่ภายในและขนานกับแกน

อินทิกรัลจากด้านภายในคือ

\[
B\ell
\]

ด้านภายนอกประมาณเป็นศูนย์ และด้านตั้งฉากกับสนามให้ดอตโปรดักต์เป็นศูนย์

จำนวนรอบที่ตัดผ่านพื้นผิวในระยะ $\ell$ คือ

\[
N_{\ell}=n\ell
\]

กระแสรวมที่ถูกล้อมคือ

\[
I_{\text{enc}}=(n\ell)I
\]

ใช้กฎของแอมแปร์

\[
B\ell
=
\mu_0(n\ell)I
\]

ตัด $\ell$

\begin{tcolorbox}[title={สนามภายในโซลินอยด์อุดมคติ}]
\[
\boxed{B=\mu_0nI}
\]

หรือ

\[
\boxed{B=\mu_0\frac{N}{L}I}
\]
\end{tcolorbox}

สูตรนี้ใช้ได้ดีใกล้กึ่งกลางของโซลินอยด์ที่ยาวมากเมื่อเทียบกับรัศมี ใกล้ปลายขดลวด สนามมีขนาดเล็กกว่าค่านี้และไม่สม่ำเสมอ

\begin{examplebox}{ตัวอย่าง: เปรียบเทียบเครื่องมือสองชนิด}
โซลินอยด์ยาว $L$ มี $N$ รอบและกระแส $I$ สนามใกล้กึ่งกลางคือ

\[
B_{\text{solenoid}}
=
\mu_0\frac{N}{L}I
\]

โทรอยด์มี $N$ รอบและพิจารณาที่รัศมี $r$ สนามคือ

\[
B_{\text{toroid}}
=
\frac{\mu_0NI}{2\pi r}
\]

ถ้านำโทรอยด์รัศมีใหญ่มากมองเพียงส่วนสั้น ๆ จะมีความยาวรอบวงประมาณ $2\pi r$ จึงมองได้ว่า

\[
n\approx\frac{N}{2\pi r}
\]

แทนใน $B=\mu_0nI$ จะได้ผลเดียวกับโทรอยด์
\end{examplebox}

% =========================================================
\section{กฎของเกาส์สำหรับสนามแม่เหล็ก}

\subsection{นิยามฟลักซ์แม่เหล็ก}

\begin{tcolorbox}[title={นิยาม: ฟลักซ์แม่เหล็กผ่านพื้นผิว}]
\[
\boxed{
\Phi_B
\equiv
\int_S\vec{B}\cdot d\vec{A}}
\]
\end{tcolorbox}

ถ้าสนามสม่ำเสมอและพื้นผิวแบน

\[
\Phi_B=BA\cos\theta
\]

โดย $\theta$ เป็นมุมระหว่าง $\vec{B}$ กับเวกเตอร์ตั้งฉากต่อพื้นผิว

หน่วยของฟลักซ์แม่เหล็กคือเวเบอร์

\[
1\,\si{Wb}=1\,\si{T.m^2}
\]

\subsection{กฎของเกาส์สำหรับสนามแม่เหล็ก}

\begin{tcolorbox}[title={กฎของเกาส์สำหรับสนามแม่เหล็ก}]
\[
\boxed{
\oint_S\vec{B}\cdot d\vec{A}=0}
\]
\end{tcolorbox}

ความหมายคือ ฟลักซ์แม่เหล็กสุทธิผ่านพื้นผิวปิดทุกพื้นผิวเป็นศูนย์

เส้นสนามแม่เหล็กจึงไม่มีจุดเริ่มหรือจุดจบ แต่เป็นเส้นปิดต่อเนื่อง ต่างจากเส้นสนามไฟฟ้าที่เริ่มจากประจุบวกและสิ้นสุดที่ประจุลบ

\begin{tcolorbox}[title={ข้อควรระวัง}]
สมการ

\[
\oint_S\vec{B}\cdot d\vec{A}=0
\]

ไม่ได้แปลว่าสนามแม่เหล็กบนพื้นผิวเป็นศูนย์ แต่แปลว่าฟลักซ์ที่ออกจากพื้นผิวรวมกับฟลักซ์ที่เข้าสู่พื้นผิวแล้วเป็นศูนย์
\end{tcolorbox}

% =========================================================
\section{แนวทางแก้โจทย์สนามแม่เหล็กจากกระแส}

\begin{enumerate}
    \item วาดทิศกระแสในแต่ละส่วนของลวดให้ครบ
    \item ระบุจุดที่ต้องการหาสนามและระยะจากลวดแต่ละส่วน
    \item ตรวจหาส่วนลวดแนวรัศมีที่ให้สนามเป็นศูนย์
    \item เลือกสูตรที่อนุมานแล้วให้ตรงกับรูปทรง
    \item ใช้กฎมือขวาหาทิศของสนามแต่ละส่วนก่อนบวกหรือลบ
    \item รวมสนามแบบเวกเตอร์ ไม่บวกเฉพาะขนาดโดยไม่ตรวจทิศ
    \item ใช้กฎของแอมแปร์เฉพาะเมื่อสมมาตรเพียงพอ
    \item ตรวจมิติของคำตอบ ซึ่งต้องอยู่ในหน่วยเทสลา
\end{enumerate}

\begin{tcolorbox}[title={ตารางเลือกสมการอย่างรวดเร็ว}]
ลวดตรงยาวมาก ระยะตั้งฉาก $r$

\[
B=\frac{\mu_0I}{2\pi r}
\]

ลวดตรงกึ่งอนันต์ที่จุดอยู่ตรงปลายฉาก

\[
B=\frac{\mu_0I}{4\pi r}
\]

ส่วนโค้งรัศมี $R$ มุม $\phi$ เรเดียน

\[
B=\frac{\mu_0I\phi}{4\pi R}
\]

วงกลมเต็มหนึ่งรอบ

\[
B=\frac{\mu_0I}{2R}
\]

วงกลมบนแกนห่างศูนย์กลาง $x$

\[
B=\frac{\mu_0IR^2}{2(R^2+x^2)^{3/2}}
\]

โซลินอยด์ยาว

\[
B=\mu_0nI
\]

โทรอยด์

\[
B=\frac{\mu_0NI}{2\pi r}
\]
\end{tcolorbox}

% =========================================================
\section{แบบฝึกหัดท้ายเนื้อหา}

\begin{exercisebox}{แบบฝึกหัดระดับพื้นฐานถึง สอวน.}
\begin{enumerate}
    \item ลวดตรงยาวมากมีกระแส $I$ ไหล จงหาสนามที่ระยะ $r$ และอธิบายทิศด้วยกฎมือขวา

    \item ลวดครึ่งวงกลมรัศมี $R$ มีกระแส $I$ ไหล จงหาสนามที่ศูนย์กลาง

    \item วงลวดรัศมี $R$ จำนวน $N$ รอบมีกระแส $I$ จงหาสนามที่ศูนย์กลาง

    \item วงลวดรัศมี $R$ มีกระแส $I$ จงหาสนามบนแกนที่ $x=R$

    \item ลวดตรงยาวมากสองเส้นห่างกัน $d$ มีกระแส $I_1$ และ $I_2$ ไหลทิศเดียวกัน จงหาแรงต่อหนึ่งหน่วยความยาวและระบุว่าเป็นแรงดูดหรือแรงผลัก

    \item ลวดทรงกระบอกรัศมี $R$ มีกระแสรวม $I$ กระจายสม่ำเสมอ จงหาสนามที่ $r=R/2$

    \item โซลินอยด์ยาว $L$ มี $N$ รอบและกระแส $I$ จงหาสนามภายใน

    \item โทรอยด์ $N$ รอบมีกระแส $I$ จงหาสนามที่รัศมี $r$ ภายในขดลวด

    \item ประจุรวม $Q$ กระจายบนวงแหวนรัศมี $R$ และหมุนด้วยอัตราเร็วเชิงมุม $\omega$ จงหากระแสสมมูลและสนามที่ศูนย์กลาง

    \item ลวดประกอบด้วยครึ่งวงกลมรัศมี $R_1$ และ $R_2$ โดยสนามที่ศูนย์กลางจากสองส่วนชี้สวนทางกัน จงหาขนาดสนามลัพธ์
\end{enumerate}
\end{exercisebox}

\subsection*{เฉลยย่อแบบฝึกหัด}

\begin{enumerate}
    \item
    \[
    B=\frac{\mu_0I}{2\pi r}
    \]

    \item
    \[
    B=\frac{\mu_0I}{4R}
    \]

    \item
    \[
    B=\frac{\mu_0NI}{2R}
    \]

    \item
    \[
    B(R)
    =
    \frac{\mu_0IR^2}{2(2R^2)^{3/2}}
    =
    \frac{\mu_0I}{2^{5/2}R}
    \]

    \item
    \[
    \frac{F}{L}=\frac{\mu_0I_1I_2}{2\pi d}
    \]
    เป็นแรงดึงดูด

    \item
    \[
    B\left(\frac{R}{2}\right)
    =
    \frac{\mu_0I(R/2)}{2\pi R^2}
    =
    \frac{\mu_0I}{4\pi R}
    \]

    \item
    \[
    B=\mu_0\frac{N}{L}I
    \]

    \item
    \[
    B=\frac{\mu_0NI}{2\pi r}
    \]

    \item
    \[
    I=\frac{Q\omega}{2\pi}
    \]
    \[
    B_O=\frac{\mu_0Q\omega}{4\pi R}
    \]

    \item
    \[
    B_{\text{net}}
    =
    \frac{\mu_0I}{4}
    \left|
    \frac{1}{R_1}-\frac{1}{R_2}
    \right|
    \]
\end{enumerate}

% =========================================================
\section{สรุปสูตรสำคัญ}

\begin{tcolorbox}[title={สรุปบท: สนามแม่เหล็กจากวงจรไฟฟ้า}]
กฎแรงแม่เหล็ก

\[
\vec{F}_B=q\vec{v}\times\vec{B}
\]

\[
d\vec{F}=I\,d\vec{\ell}\times\vec{B}
\]

กฎของบิโอต์--ซาวาร์

\[
d\vec{B}
=
\frac{\mu_0}{4\pi}
\frac{I\,d\vec{\ell}\times\hat{r}}{r^2}
\]

ลวดตรงยาวมาก

\[
B=\frac{\mu_0I}{2\pi r}
\]

ลวดตรงกึ่งอนันต์

\[
B=\frac{\mu_0I}{4\pi r}
\]

ส่วนโค้งวงกลมมุม $\phi$

\[
B=\frac{\mu_0I\phi}{4\pi R}
\]

วงกลมเต็ม

\[
B=\frac{\mu_0NI}{2R}
\]
\end{tcolorbox}

\begin{tcolorbox}[title={สรุปบท: สนามแม่เหล็กจากวงจรไฟฟ้า ต่อ}]
สนามบนแกนวงลวด

\[
B(x)
=
\frac{\mu_0NIR^2}{2(R^2+x^2)^{3/2}}
\]

วงแหวนประจุหมุน

\[
I=\frac{Q\omega}{2\pi}
\]

\[
B_O=\frac{\mu_0Q\omega}{4\pi R}
\]

แรงระหว่างลวดขนาน

\[
\frac{F}{L}
=
\frac{\mu_0I_1I_2}{2\pi d}
\]

กฎของแอมแปร์

\[
\oint\vec{B}\cdot d\vec{\ell}
=
\mu_0I_{\text{enc}}
\]

ลวดทรงกระบอก กระแสสม่ำเสมอ

\[
B(r)=\frac{\mu_0Ir}{2\pi R^2}
\qquad (r<R)
\]

\[
B(r)=\frac{\mu_0I}{2\pi r}
\qquad (r\ge R)
\]
\end{tcolorbox}

\begin{tcolorbox}[title={สรุปบท: สนามแม่เหล็กจากวงจรไฟฟ้า ต่อ}]
โทรอยด์

\[
B(r)=\frac{\mu_0NI}{2\pi r}
\]

โซลินอยด์ยาว

\[
n\equiv\frac{N}{L}
\]

\[
B=\mu_0nI
\]

นิยามฟลักซ์แม่เหล็ก

\[
\Phi_B\equiv\int\vec{B}\cdot d\vec{A}
\]

กฎของเกาส์สำหรับสนามแม่เหล็ก

\[
\oint\vec{B}\cdot d\vec{A}=0
\]
\end{tcolorbox}

% =========================================================
% The POSN archive section is deliberately placed at the end of the chapter.
% =========================================================
\section{ข้อสอบจริง สอวน. ที่ใช้ความรู้เรื่องสนามแม่เหล็กจากกระแส}


% ---------------------------------------------------------
\subsection{ปี 2562 ข้อ 11: ครึ่งวงกลมร่วมศูนย์กลางสองรัศมี}

\vspace{1em}

\IfFileExists{img/posn62q11.png}{
\begin{center}
    \includegraphics[width=1\linewidth]{img/posn62q11.png}
\end{center}
}{
\begin{tcolorbox}[title={ตำแหน่งภาพที่แนะนำ}]
ให้ตัดภาพข้อสอบปี 2562 ข้อ 11 แล้วบันทึกเป็น

\[
\texttt{img/posn62q11.png}
\]

โจทย์เป็นลวดรูปครึ่งวงกลมร่วมศูนย์กลาง รัศมี $R_1$ และ $R_2$ เชื่อมด้วยส่วนลวดตรงตามแนวรัศมี
\end{tcolorbox}
}

\begin{examplebox}{เฉลยละเอียด}
ส่วนลวดตรงตามแนวเส้นผ่านศูนย์กลางไม่สร้างสนามที่ $O$ เพราะสำหรับทุกส่วนของลวดตรง

\[
d\vec{\ell}\parallel\hat{r}
\]

จึงมี

\[
d\vec{\ell}\times\hat{r}=0
\]

เหลือเฉพาะสนามจากครึ่งวงกลมสองส่วน

สนามจากครึ่งวงกลมรัศมี $R_1$ มีขนาด

\[
B_1=\frac{\mu_0I}{4R_1}
\]

สนามจากครึ่งวงกลมรัศมี $R_2$ มีขนาด

\[
B_2=\frac{\mu_0I}{4R_2}
\]

จากทิศกระแสในรูป ครึ่งวงกลมทั้งสองวนคนละทิศ สนามจึงชี้สวนทางกัน

เพราะ

\[
R_1<R_2
\]

จึงมี

\[
B_1>B_2
\]

ขนาดสนามลัพธ์คือ

\[
B_{\text{net}}=B_1-B_2
\]

\[
\boxed{
B_{\text{net}}
=
\frac{\mu_0I}{4}
\left(
\frac{1}{R_1}-\frac{1}{R_2}
\right)}
\]

ทิศเดียวกับสนามจากครึ่งวงกลมรัศมี $R_1$ ซึ่งจากกฎมือขวาชี้เข้าหน้ากระดาษตามทิศกระแสในรูป

\[
\boxed{\text{ตอบข้อ 1}}
\]
\end{examplebox}

% ---------------------------------------------------------
\subsection{ปี 2564 ข้อ 19: สนามบนแกนของวงลวดที่สองตำแหน่ง}

\vspace{1em}

\IfFileExists{img/posn64q19.png}{
\begin{center}
    \includegraphics[width=1\linewidth]{img/posn64q19.png}
\end{center}
}{
\begin{tcolorbox}[title={ตำแหน่งภาพที่แนะนำ}]
ให้ตัดภาพข้อสอบปี 2564 ข้อ 19 แล้วบันทึกเป็น

\[
\texttt{img/posn64q19.png}
\]
\end{tcolorbox}
}

\begin{examplebox}{เฉลยละเอียด}
สนามบนแกนของวงลวดรัศมี $R$ ที่ระยะ $x$ จากศูนย์กลางคือ

\[
B(x)
=
\frac{\mu_0i}{2}
\frac{R^2}{(R^2+x^2)^{3/2}}
\]

ที่จุด $M$ ซึ่งอยู่ห่างศูนย์กลาง $a$

\[
B_M
=
\frac{\mu_0i}{2}
\frac{R^2}{(R^2+a^2)^{3/2}}
\]

ที่จุด $N$ ซึ่งอยู่ห่างศูนย์กลาง $b$

\[
B_N
=
\frac{\mu_0i}{2}
\frac{R^2}{(R^2+b^2)^{3/2}}
\]

หารสมการ

\[
\frac{B_N}{B_M}
=
\frac{
\dfrac{\mu_0iR^2}{2(R^2+b^2)^{3/2}}
}{
\dfrac{\mu_0iR^2}{2(R^2+a^2)^{3/2}}
}
\]

ตัดตัวประกอบร่วม

\[
\frac{B_N}{B_M}
=
\frac{(R^2+a^2)^{3/2}}{(R^2+b^2)^{3/2}}
\]

ดังนั้น

\[
\boxed{
B_N
=
B_M
\left(
\frac{R^2+a^2}{R^2+b^2}
\right)^{3/2}}
\]

\[
\boxed{\text{ตอบ ง}}
\]
\end{examplebox}

% ---------------------------------------------------------
\subsection{ปี 2565 ข้อ 11: ครึ่งวงกลมสองรัศมีที่สนามเสริมกัน}

\vspace{1em}

\IfFileExists{img/posn65q11.png}{
\begin{center}
    \includegraphics[width=1\linewidth]{img/posn65q11.png}
\end{center}
}{
\begin{tcolorbox}[title={ตำแหน่งภาพที่แนะนำ}]
ให้ตัดภาพข้อสอบปี 2565 ข้อ 11 แล้วบันทึกเป็น

\[
\texttt{img/posn65q11.png}
\]
\end{tcolorbox}
}

\begin{examplebox}{เฉลยละเอียด}
โจทย์ให้สนามที่ศูนย์กลางของวงกลมเต็มรัศมี $R$ ว่า

\[
B=\frac{\mu_0I}{2R}
\]

จึงเขียนได้ว่า

\[
\frac{\mu_0I}{4}=\frac{BR}{2}
\]

ในรูป ข ส่วนลวดที่เป็นแนวรัศมีไม่สร้างสนามที่ $O$ เหลือครึ่งวงกลมรัศมี $R_1$ และ $R_2$

สนามจากครึ่งวงกลมรัศมี $R_1$

\[
B_1=\frac{\mu_0I}{4R_1}
\]

สนามจากครึ่งวงกลมรัศมี $R_2$

\[
B_2=\frac{\mu_0I}{4R_2}
\]

จากทิศกระแสในรูป สนามทั้งสองชี้ทิศเดียวกัน จึงบวกกัน

\[
B_O
=
\frac{\mu_0I}{4}
\left(
\frac{1}{R_1}+\frac{1}{R_2}
\right)
\]

แทน

\[
\frac{\mu_0I}{4}=\frac{BR}{2}
\]

จะได้

\[
\boxed{
B_O
=
\frac{BR}{2}
\left(
\frac{1}{R_1}+\frac{1}{R_2}
\right)}
\]

ทิศออกจากหน้ากระดาษตามกฎมือขวาสำหรับทิศกระแสที่แสดง

\[
\boxed{\text{ตอบ ก}}
\]
\end{examplebox}

% ---------------------------------------------------------
\subsection{ปี 2565 ข้อ 20: ลวดตรงสองเส้นตามแกน $OX$ และ $OY$}

\vspace{1em}

\IfFileExists{img/posn65q20.png}{
\begin{center}
    \includegraphics[width=1\linewidth]{img/posn65q20.png}
\end{center}
}{
\begin{tcolorbox}[title={ตำแหน่งภาพที่แนะนำ}]
ให้ตัดภาพข้อสอบปี 2565 ข้อ 20 แล้วบันทึกเป็น

\[
\texttt{img/posn65q20.png}
\]
\end{tcolorbox}
}

\begin{examplebox}{เฉลยละเอียด}
จุดที่ต้องการหาสนามคือ

\[
P=(0,0,z)
\]

ระยะตั้งฉากจากจุด $P$ ถึงลวดบนแกน $OX$ และลวดบนแกน $OY$ เท่ากันคือ

\[
z
\]

สนามจากลวดแต่ละเส้นมีขนาด

\[
B_0=\frac{\mu_0I}{2\pi z}
\]

ใช้กฎมือขวา

\begin{itemize}
    \item ลวดตามแกน $+X$ สร้างสนามที่ $P$ ในทิศ $-Y$
    \item ลวดตามแกน $+Y$ สร้างสนามที่ $P$ ในทิศ $+X$
\end{itemize}

สนามทั้งสองตั้งฉากกัน จึงมีขนาดลัพธ์

\[
B_{\text{net}}
=
\sqrt{B_0^2+B_0^2}
\]

\[
B_{\text{net}}=\sqrt{2}B_0
\]

แทน $B_0$

\[
B_{\text{net}}
=
\sqrt{2}
\left(
\frac{\mu_0I}{2\pi z}
\right)
\]

\[
\boxed{
B_{\text{net}}
=
\frac{\mu_0I}{\sqrt{2}\pi z}}
\]

\[
\boxed{\text{ตอบ ค}}
\]
\end{examplebox}

% ---------------------------------------------------------
\subsection{ปี 2567 ข้อ 5: แรงโดยประมาณระหว่างวงลวดสองวง}

\vspace{1em}

\IfFileExists{img/posn67q5.png}{
\begin{center}
    \includegraphics[width=1\linewidth]{img/posn67q5.png}
\end{center}
}{
\begin{tcolorbox}[title={ตำแหน่งภาพที่แนะนำ}]
ให้ตัดภาพข้อสอบปี 2567 ข้อ 5 แล้วบันทึกเป็น

\[
\texttt{img/posn67q5.png}
\]
\end{tcolorbox}
}

\begin{examplebox}{เฉลยละเอียด}
วงลวดทั้งสองมีรัศมี $R$ และอยู่ห่างกัน $d$ โดย

\[
d\ll R
\]

ดังนั้นเมื่อพิจารณาส่วนเล็ก ๆ ที่ตำแหน่งตรงกัน ความโค้งของวงลวดมีผลน้อยมาก ส่วนลวดสองส่วนนั้นจึงประมาณเป็นลวดตรงขนานที่ห่างกัน $d$

กระแสในส่วนที่ตรงกันไหลทิศเดียวกัน เพราะกระแสในวงทั้งสองวนทิศเดียวกัน จึงเป็นแรงดึงดูด

แรงต่อหนึ่งหน่วยความยาวระหว่างส่วนลวดคือ

\[
\frac{dF}{d\ell}
\approx
\frac{\mu_0I^2}{2\pi d}
\]

แรงจากทุกส่วนชี้ตามแนวแกนเข้าหาวงอีกวง และมีค่าประมาณเท่ากันตลอดรอบวง

ความยาวรอบวงคือ

\[
L=2\pi R
\]

ดังนั้น

\[
F
\approx
\left(
\frac{\mu_0I^2}{2\pi d}
\right)
(2\pi R)
\]

\[
\boxed{
F\approx\frac{\mu_0I^2R}{d}}
\]

\[
\boxed{\text{แรงเป็นแรงดึงดูด}}
\]

คำตอบนี้ใช้ได้ภายใต้เงื่อนไข $d\ll R$ เท่านั้น
\end{examplebox}

% ---------------------------------------------------------
\subsection{ปี 2567 ข้อ 9: วงลวดสองรอบร่วมศูนย์กลาง}

\vspace{1em}

\IfFileExists{img/posn67q9.png}{
\begin{center}
    \includegraphics[width=1\linewidth]{img/posn67q9.png}
\end{center}
}{
\begin{tcolorbox}[title={ตำแหน่งภาพที่แนะนำ}]
ให้ตัดภาพข้อสอบปี 2567 ข้อ 9 แล้วบันทึกเป็น

\[
\texttt{img/posn67q9.png}
\]
\end{tcolorbox}
}

\begin{examplebox}{เฉลยละเอียด}
สนามที่ศูนย์กลางของวงลวดหนึ่งรอบรัศมี $R$ คือ

\[
B=\frac{\mu_0I}{2R}
\]

ไล่ตามเส้นลวดจากกระแสที่ไหลเข้าไปจนถึงกระแสที่ไหลออก จะพบว่ากระแสในวงรัศมี $R_1$ และ $R_2$ วนตามเข็มนาฬิกาเหมือนกันเมื่อมองตามรูป ดังนั้นสนามจากสองวงชี้เข้าหน้ากระดาษเหมือนกันและต้องบวกกัน

วงรัศมี $R_1$ ให้สนาม

\[
B_1=\frac{\mu_0I}{2R_1}
\]

วงรัศมี $R_2$ ให้สนาม

\[
B_2=\frac{\mu_0I}{2R_2}
\]

ขนาดสนามลัพธ์คือ

\[
B_O=B_1+B_2
\]

\[
\boxed{
B_O
=
\frac{\mu_0I}{2}
\left(
\frac{1}{R_1}+\frac{1}{R_2}
\right)}
\]

สนามลัพธ์ชี้เข้าหน้ากระดาษ

\[
\boxed{
\vec{B}_O\text{ ชี้เข้าหน้ากระดาษ}}
\]
\end{examplebox}

% ---------------------------------------------------------
\subsection{ปี 2568 ข้อ 5--6: วงแหวนประจุหมุนและสนามที่ศูนย์กลาง}

\vspace{1em}

\IfFileExists{img/posn68q4.png}{
\begin{center}
    \includegraphics[width=1\linewidth]{img/posn68q4.png}
\end{center}
}{
\begin{tcolorbox}[title={ภาพประกอบจากข้อ 4}]
ข้อ 4 กำหนดวงแหวนรัศมี $R$ มีประจุรวม $+Q$ กระจายสม่ำเสมอ และวางในระนาบ $XOY$

แนะนำให้ตัดภาพแล้วบันทึกเป็น

\[
\texttt{img/posn68q4.png}
\]
\end{tcolorbox}
}

\IfFileExists{img/posn68q5q6.png}{
\begin{center}
    \includegraphics[width=1\linewidth]{img/posn68q5q6.png}
\end{center}
}{
\begin{tcolorbox}[title={ตำแหน่งภาพที่แนะนำ}]
ให้ตัดภาพข้อสอบปี 2568 ข้อ 5 และข้อ 6 แล้วบันทึกเป็น

\[
\texttt{img/posn68q5q6.png}
\]
\end{tcolorbox}
}

\begin{examplebox}{เฉลยละเอียดข้อ 5}
วงแหวนหมุนด้วยอัตราเร็วเชิงมุม $\omega$

คาบการหมุนคือ

\[
T=\frac{2\pi}{\omega}
\]

ประจุรวม $Q$ ผ่านตำแหน่งใดตำแหน่งหนึ่งหนึ่งครั้งต่อหนึ่งคาบ จากนิยามกระแส

\[
I=\frac{Q}{T}
\]

แทน $T$

\[
I
=
\frac{Q}{2\pi/\omega}
\]

\[
\boxed{
I=\frac{Q\omega}{2\pi}}
\]
\end{examplebox}

\begin{examplebox}{เฉลยละเอียดข้อ 6}
สนามที่ศูนย์กลางของวงลวดรัศมี $R$ คือ

\[
B_O=\frac{\mu_0I}{2R}
\]

แทนกระแสจากข้อ 5

\[
B_O
=
\frac{\mu_0}{2R}
\left(
\frac{Q\omega}{2\pi}
\right)
\]

ดังนั้น

\[
\boxed{
B_O
=
\frac{\mu_0Q\omega}{4\pi R}}
\]

ถ้าประจุ $+Q$ หมุนทวนเข็มนาฬิกาเมื่อมองจากแกน $+Z$ กระแสตามแบบทวนเข็มนาฬิกาและสนามชี้ตามแกน $+Z$ ถ้าหมุนตามเข็ม ทิศสนามกลับกัน
\end{examplebox}

% ---------------------------------------------------------
\subsection{ปี 2568 ข้อ 25: ลวดกึ่งอนันต์สองส่วนต่อกับครึ่งวงกลม}

\vspace{1em}

\IfFileExists{img/posn68q25.png}{
\begin{center}
    \includegraphics[width=1\linewidth]{img/posn68q25.png}
\end{center}
}{
\begin{tcolorbox}[title={ตำแหน่งภาพที่แนะนำ}]
ให้ตัดภาพข้อสอบปี 2568 ข้อ 25 แล้วบันทึกเป็น

\[
\texttt{img/posn68q25.png}
\]
\end{tcolorbox}
}

\begin{examplebox}{เฉลยละเอียด}
ลวดประกอบด้วยสามส่วนที่สร้างสนาม ณ จุด $O$

\begin{enumerate}
    \item ลวดตรงกึ่งอนันต์ด้านบน $ab$ อยู่ห่าง $O$ เป็นระยะ $R$
    \item ลวดตรงกึ่งอนันต์ด้านล่าง $ed$ อยู่ห่าง $O$ เป็นระยะ $R$
    \item ครึ่งวงกลม $bcd$ รัศมี $R$
\end{enumerate}

สนามจากลวดตรงกึ่งอนันต์หนึ่งส่วนคือ

\[
B_{\text{half-line}}
=
\frac{\mu_0I}{4\pi R}
\]

จึงได้จากลวดตรงสองส่วน

\[
B_{\text{straight}}
=
2\left(
\frac{\mu_0I}{4\pi R}
\right)
\]

\[
B_{\text{straight}}
=
\frac{\mu_0I}{2\pi R}
\]

สนามจากครึ่งวงกลมคือ

\[
B_{\text{arc}}
=
\frac{\mu_0I}{4R}
\]

ใช้กฎมือขวากับทิศกระแสในรูป จะพบว่าสนามจากลวดตรงทั้งสองส่วนและครึ่งวงกลมชี้ออกจากหน้ากระดาษทั้งหมด จึงบวกขนาดได้

\[
B_O
=
\frac{\mu_0I}{2\pi R}
+
\frac{\mu_0I}{4R}
\]

ดึง $\mu_0I/R$ เป็นตัวร่วม

\[
B_O
=
\frac{\mu_0I}{R}
\left(
\frac{1}{2\pi}+\frac{1}{4}
\right)
\]

ทำส่วนให้เท่ากัน

\[
\frac{1}{2\pi}+\frac{1}{4}
=
\frac{2+\pi}{4\pi}
\]

ดังนั้น

\[
\boxed{
B_O
=
\frac{\mu_0I(\pi+2)}{4\pi R}}
\]

\[
\boxed{
\vec{B}_O\text{ ชี้ออกจากหน้ากระดาษ}}
\]
\end{examplebox}